\documentclass[../main.tex]{subfiles}
\begin{document}

\textit{The system was built with AI coding agents. This chapter describes how the agents were
organised so that they behave like a small software factory rather than a free-form assistant:
the context they are given, the stations of the production line, the guardrails that hold them to
the project's rules, the testing they must pass, and how the product is packaged, operated and
monitored.}

\section{From an assistant to a factory}
An AI coding agent left to itself has two weaknesses. It does not know the project, so it guesses;
and it does not know its limits, so it will ``fix'' a security rule as readily as a typo. Several
defects in Table~\ref{tab:defects} show both: the deploy command that would have shipped
development settings to production (F-37), and the proposed change that faked the PM role in the
browser (F-53). The design answer has four parts, each covered by a section below:
\begin{enumerate}
  \item \textbf{Context engineering}---the agent reads the right documents before acting, not after.
  \item \textbf{Factory}---the work passes through separate stations with separate permissions.
  \item \textbf{Guardrails}---the rules are executed as code, not only written as advice.
  \item \textbf{Verification}---unit, integration and end-to-end tests, and tests of the tests.
\end{enumerate}

\begin{figure}[H]
\centering
\resizebox{\textwidth}{!}{\begin{tikzpicture}[
  st/.style={draw, rounded corners, fill=blue!6, minimum width=2.15cm, minimum height=1.25cm, align=center, font=\footnotesize},
  gate/.style={draw, fill=orange!15, minimum width=13.6cm, minimum height=0.7cm, font=\footnotesize, align=center},
  ctx/.style={draw, fill=green!8, minimum width=13.6cm, minimum height=0.7cm, font=\footnotesize, align=center},
  >=Stealth]
\node[ctx] (ctx) at (6.8,2.0) {Context: CLAUDE.md + D1--D8 (requirements, modules, architecture, contracts, tests, policies, decisions, RTM)};
\node[st] (s1) at (0.0,0) {1. Require-\\ments analyst};
\node[st] (s2) at (2.27,0) {2. Architect};
\node[st] (s3) at (4.54,0) {3. Test\\engineer\\(red)};
\node[st] (s4) at (6.81,0) {4. Imple-\\menter\\(green)};
\node[st] (s5) at (9.08,0) {5. Test eng.\\(mutation)};
\node[st] (s6) at (11.35,0) {6. Reviewer};
\node[st] (s7) at (13.62,0) {7. Release\\engineer};
\foreach \a/\b in {s1/s2,s2/s3,s3/s4,s4/s5,s5/s6,s6/s7} \draw[->] (\a) -- (\b);
\draw[->, dashed] (s6.south) to[bend left=25] node[below, font=\scriptsize]{findings: back to implementer} (s4.south);
\node[gate] (g) at (6.8,-2.3) {Guardrails on every edit and command (deny / ask a human / allow) \quad + \quad Definition-of-Done gate on stop};
\foreach \s in {ctx} \draw[->] (ctx.south) -- (6.8,0.75);
\draw[->] (g.north) -- (6.8,-1.6);
\end{tikzpicture}}
\caption{The AI agent factory: context in, stations in sequence, guardrails around every step}\label{fig:factory}
\end{figure}

\section{Context engineering}
The repository contains a navigation layer written for agents first and people second. An agent
starting any task loads \code{CLAUDE.md} automatically; it lists the verification commands, the
invariants that must not be broken, and the order in which to read the eight documents:

\begin{table}[H]
\centering
\caption{The documents an agent reads, and the question each answers}\label{tab:context}
\renewcommand{\arraystretch}{1.2}
\begin{tabular}{|p{1.2cm}|p{4.2cm}|p{8.6cm}|}
\hline
\textbf{Doc} & \textbf{Content} & \textbf{Question it answers for the agent} \\ \hline
D1 & Requirements & What must the system do? (FR and NFR ids) \\ \hline
D2 & Module map & Which file implements this, and what depends on it? \\ \hline
D3 & Architecture & Which ordering or layering rule must my change keep? \\ \hline
D4 & API and data contracts & Is this a contract change? (if so: stop and escalate) \\ \hline
D5 & Tests and acceptance & What proves it works; what is the Definition of Done? \\ \hline
D6 & Risks, delegation, policies & May I do this alone, do it and report, or must I ask? \\ \hline
D7 & Development book & Was this already decided, and why? Has this bug been seen before? \\ \hline
D8 & Traceability matrix & Which requirements and tests does the file I am editing serve? \\ \hline
\end{tabular}
\end{table}

Two practices keep this context trustworthy. First, every document was written \emph{against the
code} and is updated in the same commit as the code it describes; the previous documentation had
described behaviour the code did not have, and an agent that trusts stale documents is worse than
one with none. Second, each station loads only what it needs (the analyst reads D1, D6, D8; the
architect reads D2, D3, D4, D6, D7), which keeps the context small and focused.

\section{The factory: stations and responsibilities}
Each station is a separate agent definition in \code{.claude/agents/} with its own instructions,
tool permissions and model. The command \code{/feature <description>} runs a request through all
stations in order; \code{/verify} runs every quality gate; \code{/release-check} packages and
smoke-tests the product.

\begin{longtable}{|p{2.6cm}|p{6.4cm}|p{4.9cm}|}
\caption{Factory stations}\label{tab:stations}\\
\hline
\textbf{Station} & \textbf{Produces} & \textbf{May not} \\ \hline
\endfirsthead
\hline
\textbf{Station} & \textbf{Produces} & \textbf{May not} \\ \hline
\endhead
Requirements analyst & Numbered requirements in D1, Given/When/Then acceptance criteria in D5, rows in D8. & Edit source code. \\ \hline
Architect & Impact plan: files with their autonomy zone, contract impact, invariants at risk, test list; a decision record in D7. & Edit source code; continue past a red-zone item. \\ \hline
Test engineer & Failing unit and integration tests first; after implementation, one seeded defect per test to prove it can fail. & Weaken or delete an existing assertion. \\ \hline
Implementer & The smallest change that makes the tests pass, then a production build. & Touch files outside the plan. \\ \hline
Reviewer & PASS, or numbered findings with a failure scenario: authorisation, contracts, secrets, injection, offline policy, translation, tests, docs. & Edit anything (read-only). \\ \hline
Release engineer & Gate table: unit and integration suites, guardrail tests, build, production-settings stack, smoke test, health and metrics. & Deploy or push without a human. \\ \hline
\end{longtable}

The separation is the control: a reviewer that wrote the code is not a reviewer, and an analyst
that cannot edit code cannot ``solve'' a requirement by quietly changing behaviour.

\section{Guardrails and security}
Instructions can be ignored; code cannot. The autonomy zones of D6 are therefore also implemented
as hooks that the agent runtime calls before every file edit and every shell command, and when the
agent tries to finish (Figure~\ref{fig:guard}). The policy is a set of pure functions
(\code{.claude/hooks/policy.py}); the hook script only translates between the runtime's JSON
protocol and the policy. If the script itself fails, it denies---a broken guard must not become an
open door.

\begin{figure}[H]
\centering
\begin{tikzpicture}[
  b/.style={draw, rounded corners, minimum height=0.9cm, align=center, font=\footnotesize},
  >=Stealth]
\node[b, fill=blue!6, minimum width=2.6cm] (a) at (0,0) {Agent wants to\\edit / run};
\node[b, fill=gray!10, minimum width=2.8cm] (h) at (4.0,0) {Hook\\\code{guard.py}};
\node[b, fill=gray!10, minimum width=2.8cm] (p) at (8.0,0) {Policy\\\code{policy.py}};
\node[b, fill=red!12, minimum width=2.6cm] (d) at (12.2,1.4) {DENY\\(never)};
\node[b, fill=orange!15, minimum width=2.6cm] (q) at (12.2,0) {ASK\\(human decides)};
\node[b, fill=green!12, minimum width=2.6cm] (o) at (12.2,-1.4) {ALLOW\\(+ reminder)};
\draw[->] (a) -- node[above, font=\scriptsize]{JSON} (h);
\draw[->] (h) -- (p);
\draw[->] (p.east) -- (d.west); \draw[->] (p.east) -- (q.west); \draw[->] (p.east) -- (o.west);
\end{tikzpicture}
\caption{Every edit and command passes the policy before it happens}\label{fig:guard}
\end{figure}

\begin{longtable}{|p{1.6cm}|p{7.6cm}|p{4.7cm}|}
\caption{Guardrail rules (a selection)}\label{tab:guardrules}\\
\hline
\textbf{Decision} & \textbf{Rule} & \textbf{Reason} \\ \hline
\endfirsthead
\hline
\textbf{Decision} & \textbf{Rule} & \textbf{Reason} \\ \hline
\endhead
Deny & Any edit whose content contains an API key, a private key or the previously leaked signing secret. & Secrets never enter the repository. \\ \hline
Deny & Reading or printing \code{.env}; editing \code{.env}, the frozen submission folder or generated folders. & Secrets; immutable artefacts. \\ \hline
Deny & Editing a migration that is already committed. & Applied migrations are immutable; add a new one. \\ \hline
Deny & \code{git push --force}; recursive delete of \code{/} or \code{\textasciitilde}; the old tar-over-ssh deploy. & Irrecoverable; F-37. \\ \hline
Ask & Editing authentication, the data model, API schemas, the upload allowlist, configuration, CI or the guardrails themselves. & Red zone in D6: security and contracts. \\ \hline
Ask & \code{git push}, the deploy script, database downgrade, removing a Docker volume. & Affects other people or deletes data. \\ \hline
Allow + reminder & Store, routers, translations, AI prompts, landing page. & Yellow zone: report and update docs. \\ \hline
Block stop & \code{source/} changed but no document, or backend/frontend code changed without a test, or a model changed without a migration. & Definition of Done (D5 \S4). \\ \hline
\end{longtable}

\textbf{The guard's first catch was its author.} Minutes after the hooks were enabled, they
refused a command written by the agent building this chapter: a script that \emph{wrote
documentation} mentioning the old tar-over-ssh deploy was mistaken for running it. The rule was
corrected so that text being written into a file (a ``heredoc'') is treated as data, a test was
added for exactly that case, and a second test checks that a dangerous command placed \emph{after}
such text is still caught.

\textbf{Security in the product itself.} The guardrails protect the code base; the application
has its own defences, described in Chapter~3 and tested in Chapter~4: authorisation on the server
for every project route, 404 rather than 403 for outsiders, a startup check that refuses weak
secrets, wildcard CORS or demo accounts in production, uploads restricted to an allowlist with
server-generated file names, and comments rendered as text, never as HTML.

\section{Unit and integration testing}
\begin{table}[H]
\centering
\caption{Test levels}\label{tab:levels}
\renewcommand{\arraystretch}{1.2}
\begin{tabular}{|p{2.5cm}|p{7.6cm}|p{2.0cm}|p{1.5cm}|}
\hline
\textbf{Level} & \textbf{What is exercised} & \textbf{Tool} & \textbf{Count} \\ \hline
Unit (frontend) & Pure modules---dependency sort, critical path, urgency order, keyboard dispatcher, mention parser---and the two stores in isolation. & Vitest & 442 \\ \hline
Unit (guardrails) & Every policy rule, the hook's JSON protocol, fail-closed behaviour. & pytest & 56 \\ \hline
Integration (backend) & The real routers, database session and authentication through an HTTP test client; each AI tier simulated at the HTTP boundary; migrations from an empty database. & pytest & 185 \\ \hline
Integration (UI) & Components mounted in a simulated DOM with their store. & Vitest & (in 442) \\ \hline
End-to-end & Seeded public instance, every feature through the API, AI against the real model. & script & 32 checks \\ \hline
\end{tabular}
\end{table}
Tests are themselves tested. For the guardrails, removing the force-push rule and the migration
rule made five tests fail; restoring them made all pass. Continuous integration runs three jobs on
every push---frontend, backend and guardrails---so a change that breaks any of them cannot pass
unnoticed.

\section{Packaging, operation and monitoring}
\subsection{Packaging}
The product ships as four containers described by Docker Compose: MySQL, the FastAPI backend, nginx
serving the built frontend, and \code{cloudflared} for the public tunnel. The images contain no
secret: a \code{.dockerignore} keeps \code{.env} and databases out, and the deploy script checks
the built image to confirm it. Secrets are generated per machine by \code{scripts/dev-env.sh}.

\subsection{Operation}
\code{scripts/local-prod.sh up} applies migrations as a separate step, starts the stack with the
production safety check armed, and verifies the posture (non-default secret, no demo accounts,
pinned CORS, schema at the latest revision). Docker restarts a container that stops, and log files
are capped at 5$\times$10\,MB per container.

\subsection{Monitoring}
Three signals were added for this stage:
\begin{itemize}
  \item \textbf{Health with a database check.} \code{/api/health} previously answered ``healthy''
        even when the database was down. It now runs a query and returns HTTP 503 when the database
        is unreachable; Docker's healthcheck uses it, so \code{docker ps} shows the real state.
  \item \textbf{Request id and access log.} Every response carries an \code{X-Request-ID}; one log
        line per request records method, route, status and duration. The route \emph{template} is
        logged (\code{/api/tasks/\{task\_id\}}), so no id, token or query string reaches a log.
  \item \textbf{Metrics.} \code{/api/metrics} publishes, in the Prometheus text format, request
        counts by route and status, a latency histogram and \textbf{AI answers by fallback tier}.
        The last one matters most: when the AI provider fails, every endpoint still answers (by
        design), so only this counter reveals that users are receiving the offline fallback instead
        of the model. Metrics are hidden (404) in production unless a secret token is presented.
\end{itemize}
\code{scripts/monitor.sh} combines them into a snapshot and prints a warning when server errors
appear or when more than half of the AI answers come from the fallback. Its output on the public
instance after the end-to-end test:
\begin{verbatim}
health      HTTP 200  (0.0026s)
container   koshi-backend-prod-local   Up (healthy)
container   koshi-db-prod-local        Up 2 days (healthy)
requests    48 total, 0 server errors
AI answers  5, 0 from the offline fallback (0%)
\end{verbatim}
Eight integration tests cover this surface, including that the metrics endpoint is hidden without
the token, that a hostile request id is replaced, and that a task id never appears in a metric label.

\section{Assessment}
The factory does not make the agents infallible. The stop gate verifies that documents and tests
\emph{changed}, not that they are correct; that judgement belongs to the reviewer station and,
finally, to people. The hooks bind the coding agent only, which is why continuous integration
re-runs every suite independently. Within those limits, the rules that previously depended on an
agent choosing to read a policy document are now enforced on every action, and each rule is
covered by a test that has been shown to fail when the rule is removed.

\end{document}
